\section{MiMo-V2：多模型编排与生态伏击}

罗福莉把 MiMo-V2 系列称为一次“觉醒”和“伏击”。一年前做多模态和语音时，它们像独立模型；看到 OpenClaw 后，她突然意识到这些模型可以在 Agent 框架中被组织和编排：语言模型负责规划和推理，多模态模型负责感知，语音模型负责交互，小模型负责低成本高频调用。

\begin{figure}[H]
\centering
\includegraphics[width=0.92\textwidth]{mimo-orchestration.png}
\caption{MiMo-V2 生态在 Agent 任务中的多模型编排。}
\end{figure}

\begin{knowledgebox}{读图：为什么不是所有能力都合进一个模型}
图中 Agent scaffold 位于中心，周围是语言、多模态、语音和小模型。它支持的结论是：端到端生产力不一定来自一个最大模型，而来自按任务阶段调度不同模型，以获得更好的速度、成本和能力组合。
\end{knowledgebox}

\subsection{成本、速度与价格}

罗福莉多次强调，生产力革命必须在意端到端完成率与成本效率。语音生成没有必要用一体化大模型，多模态理解是否值得更大模型也要打问号。Agent 框架让不那么顶尖但更便宜、更快的小模型获得更大发挥空间，因为框架能弥补短板并稳定输出。

\begin{warningbox}{不要只问“哪个模型最强”}
Agent 系统更应该问：哪个模型在这个环节最合适。规划、感知、语音、验证、执行、总结可能需要不同模型。最强模型若太慢太贵，未必是生产系统最优解。
\end{warningbox}

\subsection{同生态模型的协同}

罗福莉提到，同一个生态训练出来的模型共享背景知识和能力边界，因此可以更放心地在 Agent 框架中分工。这里的关键不是品牌协同，而是系统知道每个模型擅长什么、成本如何、输出如何接入下一步。

\begin{importantbox}{MiMo-V2 的工程含义}
MiMo-V2 在访谈中的意义不只是发布多个模型，而是说明 Agent 时代模型家族需要被编排成系统。模型、工具、任务和成本路由共同构成产品能力。
\end{importantbox}

\subsection{本章小结}

MiMo-V2 章节说明，多模型生态在 Agent 时代会从“并列发布”变成“任务编排”。未来竞争不只是谁有一个最强模型，也是谁能把多个模型组织成低成本、高完成率的 Agent 系统。

